\chapter{AGENTIC VISION-GUIDED DRONE CONTROL}
\label{chap:4th chap}

% ─────────────────────────────────────────────────────────────────────────────
\section{Introduction}
\label{sec:ch4_intro}
% ─────────────────────────────────────────────────────────────────────────────

Chapters~\ref{chap:2nd chap} and~\ref{chap:3rd chap} established two complementary but
separate capabilities.  Chapter~\ref{chap:2nd chap} demonstrated that large language models
(LLMs) can \emph{verbalize} a drone's camera feed --- translating raw image pixels into
structured natural-language descriptions of scene hazards, object positions, and recommended
pilot actions --- across 800 trials with a best-case safety rate of 100\% and zero dangerous
outputs~\cite{vemprala2023chatgpt}.  Chapter~\ref{chap:3rd chap} demonstrated that the same
family of models can \emph{execute} high-level natural-language commands by issuing a
sequence of typed tool calls to a drone flight controller, achieving precise altitude and
position control without a human pilot in the loop~\cite{yao2022react}.

Both capabilities, however, remained isolated.  In Chapter~\ref{chap:2nd chap} the LLM
produced \emph{suggestions}; a human pilot decided whether to act on them.  In
Chapter~\ref{chap:3rd chap} the LLM \emph{acted}, but a human operator still provided the
visual situational awareness.  The natural next question is: can a single LLM orchestrator
both \emph{perceive} --- by calling an image analysis tool --- and \emph{act} --- by calling
drone control tools --- within a unified, multi-turn agentic loop?  If so, a quadrotor could
conduct an indoor patrol mission from arm-up to final report without any human in the
perception-action pipeline, while still accepting human oversight when safety is at stake.

This chapter describes the \textbf{ND series} of experiments (Neural Drone: Natural-language
Driven), which combine scene verbalization with tool-call execution into an integrated
autonomous patrol system.  The experiments proceed in three stages:

\begin{enumerate}
  \item \textbf{ND1} validates that \texttt{analyze\_scene} --- the image verbalization call
    from Chapter~\ref{chap:2nd chap} --- can function as a first-class tool call within the
    orchestrator's tool namespace, with the same 100\% invocation reliability observed in
    hand-authored verbalization sessions.

  \item \textbf{ND2} extends ND1 to full multi-turn patrol missions, revealing a critical
    failure mode: GPT-4o-mini enters an infinite semantic loop when the emergency monitor
    injects repeated hazard advisories, a manifestation of the \emph{Lost in the Middle}
    phenomenon~\cite{liu2023lostinthemiddle}.  Six targeted mitigations (M1--M6) are
    developed and validated.

  \item \textbf{ND3} introduces a \emph{human-in-the-loop} (HITL) approval gate that
    intercepts every drone-control tool call for operator sign-off, converting the fully
    autonomous system into a supervised one.  Two critical API bugs are discovered and fixed,
    and the M1--M6 mitigations are shown to be necessary even when a HITL gate is present.
\end{enumerate}

The unified architecture achieves patrol completion rates of 60--100\%, mission report
quality up to 4.60/5, and per-run costs ranging from \$0.016 (GPT-4o-mini with fixes) to
\$1.91 (Claude), providing a concrete cost-quality tradeoff map for deployment decisions.

% ─────────────────────────────────────────────────────────────────────────────
\section{Problem Statement}
\label{sec:ch4_problem}
% ─────────────────────────────────────────────────────────────────────────────

\subsection{The Perception--Action Integration Gap}
\label{sec:ch4_problem_gap}

Chapter~\ref{chap:2nd chap} treated image verbalization as a \emph{terminal} operation:
given an image, the LLM returned a structured analysis, and the pipeline ended.  The output
was textual and consumed by a human.  Chapter~\ref{chap:3rd chap} treated commands as
\emph{inputs}: a human operator supplied a natural-language objective, and the LLM translated
it into a sequence of tool calls.  Neither chapter addressed the full autonomous loop in
which the drone's \emph{own} camera output feeds back into the LLM's decision at the next
time step.

Closing this loop requires \texttt{analyze\_scene} to become a callable tool, not merely an
external service.  When the LLM can both request a scene analysis and act on its result
within the same reasoning turn, the system becomes capable of perception-contingent
planning --- for example, navigating to a waypoint, calling \texttt{analyze\_scene}, deciding
whether to escalate or continue, and planning the next waypoint based on what was found.

This capability was previously demonstrated in structured form by \citet{huang2022innermonologue}
(Inner Monologue), where perception observations were injected into the LLM's reasoning
context at each step.  The ND series adopts this pattern but moves from scripted injection
to autonomous tool-call invocation, giving the orchestrator \emph{agency} over when and
where to look.

\subsection{Why Human-in-the-Loop Oversight Matters}
\label{sec:ch4_problem_hitl}

Fully autonomous drone operation in indoor environments raises safety and accountability
concerns that go beyond what a software guardrail can fully address.  The guardrail system
introduced in Chapter~\ref{chap:3rd chap} enforced hard bounds on altitude and position
but operated \emph{after} the tool call was proposed.  For safety-critical operations ---
particularly those involving persons detected in the scene --- a human operator may need to
review and approve each drone-control action before execution.

\citet{sanneman2022safeai} formalise this requirement in the SAFE-AI framework, arguing that
human feedback is necessary for maintaining situational awareness in human-AI teaming,
especially when the AI system's internal reasoning is not directly observable.
\citet{chen2018humanperformance} further show that operator response time and approval latency
are tractable: skilled operators can evaluate and approve or deny an action within 1--3
seconds, well within the timescale of low-frequency LLM control loops operating at
0.1--0.4~Hz.

The ND3 experiments therefore introduce a \texttt{approve\_callback} gate: every tool call
in the drone-control namespace (\texttt{DRONE\_CONTROL\_TOOLS}) is intercepted and presented
to a simulated operator (auto-approver with safety bounds) before execution.
\texttt{analyze\_scene} and sensor-read calls are exempt --- only actuator commands require
approval.

\subsection{Research Questions and Scope}
\label{sec:ch4_problem_rq}

The ND series addresses four research questions:

\begin{description}
  \item[RQ1] Can an LLM orchestrator autonomously invoke \texttt{analyze\_scene} as a tool
    call with the same reliability observed in human-initiated verbalization sessions?

  \item[RQ2] Do multi-turn patrol missions with repeated emergency monitor advisories degrade
    orchestrator performance, and does the degradation follow the Lost-in-the-Middle failure
    mode predicted by \citet{liu2023lostinthemiddle}?

  \item[RQ3] Can the M1--M6 semantic-loop mitigations developed for the standalone (no-HITL)
    setting also resolve loop failures when an operator approval gate is present?

  \item[RQ4] How does the HITL gate affect per-run cost, mission quality, and safety in an
    agentic drone patrol system, and which orchestrator offers the best cost-quality tradeoff?
\end{description}

All experiments run in the same custom Python simulator used in Chapters~\ref{chap:2nd chap}
and~\ref{chap:3rd chap}.  The drone model is a custom 50-gram quadrotor with a bespoke
flight controller; no standard autopilot stack (PX4, ArduPilot) is involved.

% ─────────────────────────────────────────────────────────────────────────────
\section{Proposed Approach: Unified Perception--Action Architecture}
\label{sec:ch4_proposed}
% ─────────────────────────────────────────────────────────────────────────────

\subsection{ND Series Architecture}
\label{sec:ch4_proposed_arch}

Figure~\ref{fig:ch4_nd_architecture} illustrates the ND series architecture.  The
orchestrator LLM sits at the centre of a closed perception-action loop.  On each turn the
model receives a \emph{mission state context} (current position, turn count, previous tool
results, emergency advisories) and produces either a set of tool calls or a final text
report.  The tool set is partitioned into three groups:

\begin{enumerate}
  \item \textbf{Perception tools}: \texttt{analyze\_scene(image\_path)} --- triggers the
    full Chapter~\ref{chap:2nd chap} vision pipeline (MediaPipe
    EfficientDet-Lite0~\cite{lugaresi2019mediapipe}, YOLO-World~\cite{cheng2024yoloworld},
    DepthAnything v2~\cite{yang2024depthanythingv2}) and returns a structured
    risk assessment; \texttt{get\_sensor\_status()} --- returns current altitude, tilt, and
    battery state.  Both run without HITL approval.

  \item \textbf{Control tools} (DRONE\_CONTROL\_TOOLS, HITL-gated in ND3):
    \texttt{arm\_motors}, \texttt{find\_hover\_throttle}, \texttt{enable\_altitude\_hold},
    \texttt{set\_altitude\_target}, \texttt{navigate\_to\_waypoint}, \texttt{hover},
    \texttt{land}, \texttt{land\_and\_disarm}.

  \item \textbf{Reporting}: the model writes a natural-language \texttt{final\_report} in
    its last text turn.  Report content, word count, risk label, and adherence to mission
    objectives are evaluated as the primary quality metric.
\end{enumerate}

The architecture draws on the ReAct pattern of~\citet{yao2022react}: reason (mission state
context), act (tool call), observe (tool result), then reason again.  GSCE prompt engineering
guidelines~\cite{wang2025gsce} inform the mission prompt structure: \emph{Guidelines},
\emph{Skill APIs}, \emph{Constraints}, and \emph{Examples} are provided in that order.

\begin{figure}[ht]
  \centering
  \includegraphics[width=0.90\linewidth]{figures/fig4_1_nd_architecture.pdf}
  \caption{ND series unified perception--action architecture.  The LLM orchestrator calls
    \texttt{analyze\_scene} for visual intelligence and drone-control tools for physical
    actuation.  In ND3, all control tools pass through a HITL approval gate before
    execution.  A background emergency monitor injects hazard advisories into the next
    turn's context.}
  \label{fig:ch4_nd_architecture}
\end{figure}

\subsection{Tool Namespace Extension: \texttt{analyze\_scene} as a Tool}
\label{sec:ch4_proposed_tool}

The key architectural change from Chapter~\ref{chap:3rd chap} is registering
\texttt{analyze\_scene} in the same function-calling namespace as the drone-control tools.
In Chapter~\ref{chap:2nd chap}, image analysis was triggered externally by a human or a
script; the LLM received the result as a prompt injection.  In the ND series the LLM
\emph{decides} when to call \texttt{analyze\_scene} and with which image file.

This mirrors the ToolEQA pattern of~\citet{embodiedqa2025}, where grounded tool outputs
provide verifiable reasoning steps for embodied question answering.  Each
\texttt{analyze\_scene} call returns a JSON object containing: MediaPipe-detected person
confidence and estimated distance, YOLO-World top labels, DepthAnything v2 relative depth
map statistics, and an overall risk label (\texttt{CLEAR} / \texttt{ALERT} /
\texttt{EMERGENCY}).  The orchestrator uses this structured output to decide the next
action.

\subsection{Emergency Monitor and Advisory Injection}
\label{sec:ch4_proposed_monitor}

A background thread polls the MediaPipe EfficientDet-Lite0 detector at 1~Hz, independently
of the LLM loop.  When a person is detected above a confidence threshold in the current
scene, the monitor generates a hazard advisory string that is prepended to the next LLM
turn's user message.  This simulates real-time sensor interrupts that may arrive at any
point during a patrol mission.

In ND1 the monitor is passive (no alert scenario).  In ND2 and ND3 the alert\_room scenario
activates the monitor, which fires on every turn --- producing the repeated advisory
injection that triggers the Lost-in-the-Middle failure mode in Section~\ref{sec:ch4_nd2_litm}.

\subsection{Human-in-the-Loop Gate (ND3 Only)}
\label{sec:ch4_proposed_hitl}

The \texttt{approve\_callback} function is passed to the orchestrator loop before ND3 runs.
On each proposed tool call the gate:
\begin{enumerate}
  \item Checks whether the tool is in \texttt{DRONE\_CONTROL\_TOOLS}.
  \item Checks whether the proposed parameters fall within safety bounds
    ($\pm$2.5~m patrol zone, 0.5--2.0~m altitude, $\leq$3 hover calls, $\leq$2.0~m
    move step).
  \item Returns \texttt{APPROVED} with a 1.22~s mean simulated decision latency (uniform
    $U[0.5, 2.0]$~s, modelling realistic human response~\cite{chen2018humanperformance})
    or \texttt{DENIED} with an explanatory message injected into the tool result.
\end{enumerate}
\texttt{analyze\_scene} and \texttt{get\_sensor\_status} are always exempt from the gate,
so perception can occur without operator bottleneck.

% ─────────────────────────────────────────────────────────────────────────────
\section{Experimental Setup}
\label{sec:ch4_setup}
% ─────────────────────────────────────────────────────────────────────────────

\paragraph{Orchestrators.}
Four commercial LLMs are evaluated throughout the ND series: Claude~3.5~Sonnet
(\texttt{claude-3-5-sonnet-20241022}), GPT-4o (\texttt{gpt-4o-2024-11-20})~\cite{achiam2023gpt4},
GPT-4o-mini (\texttt{gpt-4o-mini})~\cite{achiam2023gpt4}, and Gemini~2.5~Flash
(\texttt{gemini-2.5-flash-preview-04-17})~\cite{reid2024gemini}.
All run at temperature~$t=0.0$ (greedy decoding), consistent with the V8 finding in
Chapter~\ref{chap:2nd chap}.

\paragraph{Scenarios.}
ND1 uses eight static scene images from the Chapter~\ref{chap:2nd chap} validation set.
ND2 and ND3 use two patrol scenarios:
\begin{itemize}
  \item \textbf{clear\_room}: scene image \texttt{door\_open\_run03\_real.jpg},
    ground-truth label \texttt{CLEAR}.  No emergency monitor activation.
  \item \textbf{alert\_room}: scene image \texttt{person\_detected.jpg},
    ground-truth label \texttt{ALERT}.  Emergency monitor fires on every LLM turn,
    injecting a hazard advisory.
\end{itemize}

\paragraph{Runs and cost accounting.}
Each condition is replicated five times ($N=5$) for reliability estimation following
HELM~\cite{liang2023helm}.  API costs are tracked at the token level; all figures quoted
are actual billed USD.  ND1 total: \$2.83 (160~trials).  ND2 total: \$17.52 (40~trials).
ND3/ND3B/ND3C combined: \$33.01 (total across all sub-runs and re-runs).

\paragraph{Metrics.}
\begin{itemize}
  \item \textbf{Patrol rate}: fraction of runs in which the drone completed all designated
    waypoints and produced a final report.
  \item \textbf{Quality (0--5)}: human-rated report quality assessing completeness,
    risk label accuracy, and actionability.
  \item \textbf{s3\_risk}: binary indicator (1 if risk label matches ground truth, else 0).
  \item \textbf{Cost/run}: billed USD per completed (or terminated) run.
  \item \textbf{Turns/run}: number of LLM API calls per run.
  \item \textbf{Flip rate}: fraction of repeated runs that switch risk label
    (non-determinism measure, ND1 only).
\end{itemize}

\paragraph{Simulator and hardware.}
The same custom Python simulator from Chapters~\ref{chap:2nd chap} and~\ref{chap:3rd chap}
is used.  The four-tier control architecture (Tier~1 PID at 4~kHz, Tier~1.5 MediaPipe at
60~Hz, Tier~2 YOLO+DepthAnything at 4~Hz, Tier~3 LLM at 0.1--0.4~Hz) remains unchanged.

% ─────────────────────────────────────────────────────────────────────────────
\section{ND1: Camera-as-Tool Validation}
\label{sec:ch4_nd1}
% ─────────────────────────────────────────────────────────────────────────────

\subsection{Objective and Design}
\label{sec:ch4_nd1_design}

ND1 asks the simplest possible question: when \texttt{analyze\_scene} is listed in the
orchestrator's tool namespace alongside the drone-control tools, will the model actually call
it?  The concern is that a model focused on \emph{acting} (moving, hovering, landing) may
treat the perception tool as optional and skip it, yielding physically correct flight but
semantically empty patrol.

The experiment fixes 8 static scene images (4 clear, 4 alert), asks each of the 4
orchestrators to complete a single-scene patrol (arm, navigate, analyse, land, report), and
repeats each image--orchestrator combination 5 times.  Total: $4 \times 8 \times 5 = 160$
trials.

\subsection{Results}
\label{sec:ch4_nd1_results}

\begin{table}[ht]
  \caption{ND1 results: camera-as-tool validation (160 trials, 8 scenes, 4 orchestrators,
    5 runs each).  All orchestrators achieved 100\% \texttt{analyze\_scene} invocation rate.
    GPT-4o-mini leads on quality; Gemini leads on cost efficiency (146$\times$ cheaper than
    Claude).  Claude's lower quality reflects report verbosity exceeding the 150-word
    evaluation ceiling.}
  \label{tab:ch4_nd1}
  \centering
  \small
  \begin{tabular}{lcccc}
    \toprule
    Orchestrator & Quality/5 & Flip Rate & Cost/Run & Tool Call Rate \\
    \midrule
    GPT-4o-mini (best quality) & \textbf{4.42} & \textbf{12.5\%} & \$0.00158 & 100\% \\
    Gemini 2.5 Flash            & 4.40          & 18.8\%          & \textbf{\$0.00043} & 100\% \\
    GPT-4o                      & 4.38          & 18.8\%          & \$0.02452 & 100\% \\
    Claude 3.5 Sonnet           & 3.00          & 55.2\%          & \$0.04425 & 100\% \\
    \bottomrule
  \end{tabular}
\end{table}

Table~\ref{tab:ch4_nd1} summarises ND1 results.  Three findings stand out:

\paragraph{Finding ND1-1 --- 100\% tool invocation rate.}
All four orchestrators called \texttt{analyze\_scene} on every single trial.  No model
skipped or deferred the perception step.  This answers RQ1 affirmatively: registering
\texttt{analyze\_scene} in the tool namespace is sufficient to make perception a natural
part of the model's action plan, without any explicit instruction to ``call the scene
analysis tool.''  This is consistent with the SayCan affordance model of~\citet{ahn2022saycan},
where available tools become action candidates that the model selects based on task-relevance.

\paragraph{Finding ND1-2 --- Quality-cost tension across models.}
GPT-4o-mini achieves the highest report quality (4.42/5) at \$0.00158/run --- roughly
10~cents per hundred runs.  Gemini 2.5 Flash matches quality at 4.40/5 but costs only
\$0.00043/run (146$\times$ cheaper than Claude at \$0.04425/run).  GPT-4o sits in the
mid-quality, mid-cost tier.  Claude scores 3.00/5 despite the highest cost, because its
verbose reports (over-explained reasoning chains exceeding 150~words) dilute the structured
risk assessment the quality rubric rewards.  This verbosity failure for Claude was also
observed in the G-series experiments of Chapter~\ref{chap:2nd chap}.

\paragraph{Finding ND1-3 --- Flip rate reveals non-determinism at $t=0.0$.}
Even with greedy decoding ($t=0.0$), Claude produces different risk labels across repeated
runs on identical inputs (55.2\% flip rate).  This confirms the V8 finding: Claude's
sampling mechanism does not guarantee full determinism at zero temperature, possibly due to
floating-point non-reproducibility across API calls~\cite{achiam2023gpt4}.  GPT-4o-mini's
12.5\% flip rate is the most stable, consistent with its deterministic clear/alert
behaviour observed subsequently in ND2.

\begin{figure}[ht]
  \centering
  \includegraphics[width=0.80\linewidth]{figures/fig4_2_nd1_quality_cost.pdf}
  \caption{ND1 quality vs.\ cost per run (bubble size proportional to flip rate).
    GPT-4o-mini offers the best quality-cost ratio; Gemini the lowest cost overall.
    Claude's high cost and lower quality reflect report verbosity.
    All four orchestrators achieved 100\% \texttt{analyze\_scene} invocation.}
  \label{fig:ch4_nd1_qc}
\end{figure}

\subsection{ND1 Summary}
\label{sec:ch4_nd1_summary}

ND1 validates \texttt{analyze\_scene} as a first-class agentic tool (RQ1: \textbf{yes}).
The simplest single-scene missions reveal a clear ordering: GPT-4o-mini and Gemini best
balance quality and cost; Claude is the weakest single-scene performer despite its
capability as an NL-command executor in Chapter~\ref{chap:3rd chap}.  With the per-scene
baseline established, ND2 extends to full multi-turn patrol.

% ─────────────────────────────────────────────────────────────────────────────
\section{ND2: Multi-Turn Agentic Patrol}
\label{sec:ch4_nd2}
% ─────────────────────────────────────────────────────────────────────────────

\subsection{Objective and Mission Design}
\label{sec:ch4_nd2_design}

ND2 scales from single-scene snapshots to complete patrol missions.  Each run begins with
the drone disarmed on the ground and ends with the drone landed and a final security report
written.  The mission involves: arming, altitude hold activation, climbing to 1.0~m,
navigating a three-waypoint route, calling \texttt{analyze\_scene} at each waypoint, and
producing a consolidated room-status report.  A 1~Hz background emergency monitor fires
throughout the alert\_room scenario.

Parameters: $\max\_turns = 50$, $t = 0.0$, 4 orchestrators $\times$ 2 scenarios $\times$
5 runs $= 40$~trials, total cost \$17.52.

\subsection{Per-Orchestrator Results}
\label{sec:ch4_nd2_results}

\begin{table}[ht]
  \caption{ND2 multi-turn patrol results (40 trials, 4 orchestrators, 2 scenarios,
    5 runs each).  ``Patrol\%'' is the fraction of runs completing all waypoints plus a
    final report.  GPT-4o delivers the best overall balance; GPT-4o-mini's alert\_room
    LITM failure and Gemini's clear\_room variability are the main failure modes.}
  \label{tab:ch4_nd2}
  \centering
  \small
  \begin{tabular}{llccc}
    \toprule
    Orchestrator & Scenario & Patrol\% & Correct\% & Cost/Run \\
    \midrule
    \multirow{2}{*}{Claude 3.5 Sonnet}
      & clear\_room & 0\%   & 60\%  & --- \\
      & alert\_room & 100\% & 100\% & --- \\
    \midrule
    \multirow{2}{*}{GPT-4o}
      & clear\_room & \textbf{100\%} & \textbf{100\%} & --- \\
      & alert\_room & 40\%  & 80\%  & --- \\
    \midrule
    \multirow{2}{*}{GPT-4o-mini}
      & clear\_room & \textbf{100\%} & \textbf{100\%} & \$0.024 \\
      & alert\_room & 0\%   & 100\% & \$0.615 \\
    \midrule
    \multirow{2}{*}{Gemini 2.5 Flash}
      & clear\_room & 60\%  & 60\%  & \textbf{\$0.024} \\
      & alert\_room & 40\%  & 100\% & \textbf{\$0.024} \\
    \bottomrule
  \end{tabular}
\end{table}

Table~\ref{tab:ch4_nd2} reveals three behavioural patterns.

\paragraph{Finding ND2-1 --- GPT-4o delivers the best overall balance.}
GPT-4o completes the clear\_room patrol with 100\% rate and correct risk labelling.  In
alert\_room, 40\% of runs halt after detecting the person hazard (runs completed in 3--5
turns, writing an emergency report and requesting return-to-home without completing all
waypoints).  This is \emph{correct safety doctrine}: abort patrol on confirmed hazard,
escalate immediately.  The remaining 60\% complete the patrol and then report.  GPT-4o's
combined behaviour (maximise mission when safe, escalate when not) makes it the most
balanced orchestrator for patrol-plus-safety objectives.

\paragraph{Finding ND2-2 --- Claude loops in clear\_room.}
Despite its strong NL-command performance in Chapter~\ref{chap:3rd chap}, Claude fails
clear\_room completely (0\% patrol): it enters a reasoning loop, repeatedly re-querying
conditions without making forward progress, and exhausts its turn budget.  Alert\_room
succeeds (100\%) because the hazard signal provides a clear decision point that terminates
the loop.  This \emph{decision-signal dependence} is not predicted by single-scene ND1
results and illustrates how multi-turn mission complexity changes model behaviour.

\paragraph{Finding ND2-3 --- Gemini is cheapest but variable.}
Gemini 2.5 Flash completes 60\% of clear\_room patrols and 40\% of alert\_room patrols,
both at \$0.024/run --- the lowest cost across all ND2 conditions.  The variability likely
reflects probabilistic mission termination decisions; some runs write the report early and
land, others overrun.  Gemini's alert\_room 40\% patrol rate also includes hazard-abort
behaviour (correct) in the remaining 60\%.

\begin{figure}[ht]
  \centering
  \includegraphics[width=0.80\linewidth]{figures/fig4_3_nd2_patrol_heatmap.pdf}
  \caption{ND2 patrol completion rates by orchestrator and scenario.  Colour encodes
    patrol\% (green = 100\%, red = 0\%).  GPT-4o clears both scenarios; GPT-4o-mini has a
    stark clear/alert split; Claude and Gemini are partially effective.}
  \label{fig:ch4_nd2_heatmap}
\end{figure}

\subsection{Lost in the Middle: GPT-4o-mini Alert\_room Failure}
\label{sec:ch4_nd2_litm}

GPT-4o-mini achieves 100\% clear\_room patrol but fails alert\_room completely
(0\% patrol, 0\% report written) despite 100\% correct risk identification.  The failure
mode is a semantic infinite loop.

\paragraph{Mechanism.}
The emergency monitor fires every turn, injecting an advisory such as
``\texttt{[EMERGENCY MONITOR] Person detected. Take evasive action.}''  On each turn the
model processes the advisory, determines it should hover to assess the situation, hovers,
receives the advisory again on the next turn, and hovers again.  The mission never
progresses beyond the hovering phase.

This is a direct instance of the \emph{Lost in the Middle} phenomenon documented
by~\citet{liu2023lostinthemiddle}: as the context window fills with repeated tool-call
results and emergency advisories (each turn appending $\approx$8,700~tokens), the model's
effective attention to the original mission objective --- ``complete the patrol, write a
report, land'' --- degrades.  \citet{liu2023lostinthemiddle} show a U-shaped attention
curve where information placed in the middle of a long context receives up to 30\% less
attention than information at the start or end.  With repeated injections growing the context
each turn, the mission exit condition is progressively buried.

\begin{figure}[ht]
  \centering
  \includegraphics[width=0.80\linewidth]{figures/fig4_4_nd2_litm_growth.pdf}
  \caption{Context token growth across turns for GPT-4o-mini alert\_room (ND2, no fixes).
    Each tool call appends $\approx$8,700~tokens; by turn~50 the context reaches
    440,057~tokens --- 34.2$\times$ growth from turn~1.  The mission objective (written at
    turn~1) is in the middle of this growing context, receiving degraded attention
    per \citet{liu2023lostinthemiddle}.}
  \label{fig:ch4_nd2_litm}
\end{figure}

Figure~\ref{fig:ch4_nd2_litm} shows the 34.2$\times$ context growth from turn~1
(12,871~tokens) to turn~50 (440,057~tokens).  The total 5-run cost is \$3.075 --- entirely
wasted, producing zero patrol completions.  The LITM failure is a structural property: all
5 alert\_room runs produce identical sequences (same tokens, same loop pattern), confirming
the outcome is deterministic given greedy decoding.

\subsection{M1--M6 Mitigations}
\label{sec:ch4_nd2_mitigations}

Six targeted mitigations are developed and applied to the GPT-4o-mini alert\_room failure.
Table~\ref{tab:ch4_nd2_mitigations} summarises the mitigations; their combined effect
reverses the failure completely.

\begin{table}[ht]
  \caption{M1--M6 semantic loop mitigations: mechanism, parameter, and theoretical basis.
    All six are applied together in the ND2 fix experiment and subsequently in ND3C.}
  \label{tab:ch4_nd2_mitigations}
  \centering
  \small
  \renewcommand{\arraystretch}{1.25}
  \begin{tabular}{cllp{4.2cm}}
    \toprule
    \# & Name & Parameter & Mechanism \\
    \midrule
    M1 & Advisory rate-limit & \texttt{MIN\_ADVISORY\_GAP\_TURNS = 5}
       & Suppress re-injection if fewer than 5 turns since last advisory;
         reduces context noise~\cite{liu2023lostinthemiddle} \\
    M2 & Structured memory injection & Every turn, primacy position
       & Prepend \texttt{\_build\_mission\_state()} to user message;
         U-shaped attention curve~\cite{liu2023lostinthemiddle} \\
    M3 & \texttt{final\_text} overwrite fix & Guard clause
       & Prevent non-empty report text being overwritten by later tool turns \\
    M4 & Loop detection + nudge & \texttt{SCENE\_CALL\_NUDGE\_THRESHOLD = 3}
       & Inject \emph{``LOOP DETECTED --- write report NOW''} once
         \texttt{analyze\_scene} is called $\geq$3$\times$~\cite{runaway2025} \\
    M5 & Turn budget warning & \texttt{TURN\_BUDGET\_WARN\_AT = 40}
       & Inject \emph{``Only N turns remaining''} once turn $\geq$40;
         creates deadline urgency \\
    M6 & Explicit quit condition & Always present in mission state
       & EXIT RULE at primacy: \emph{``analyze\_scene $\geq$3$\times$ $\rightarrow$
         STOP tools, write report, land''}~\cite{yao2022react} \\
    \bottomrule
  \end{tabular}
\end{table}

The theoretical basis for M1--M2 is the primacy effect of the U-shaped attention
curve~\cite{liu2023lostinthemiddle,foundinthemiddle2024}: placing the exit condition at the
\emph{start} (primacy position) of every turn's user message ensures it receives maximum
attention regardless of context length.  M4 and M6 together give the model an explicit,
self-applicable exit condition that it can trigger autonomously --- addressing the core gap
that \citet{runaway2025} identify in embodied LLM agents lacking intrinsic termination
criteria.

\begin{table}[ht]
  \caption{M1--M6 fix results on GPT-4o-mini alert\_room (ND2, 5 runs).  The mitigation
    package transforms complete failure into full success, with a 96\% cost reduction.
    Context growth drops from 34.2$\times$ to 1.4$\times$.}
  \label{tab:ch4_nd2_fix}
  \centering
  \small
  \begin{tabular}{lcc}
    \toprule
    Metric & Before (ND2, no fix) & After (M1--M6 applied) \\
    \midrule
    Patrol rate     & 0\%      & \textbf{100\%} \\
    Quality/5       & 0.0      & \textbf{4.0}   \\
    Words written   & 0        & \textbf{$\sim$130} \\
    Turns/run       & 50 (max) & \textbf{$<$15} \\
    Cost/run        & \$0.615  & \textbf{\$0.027} \\
    Context growth  & 34.2$\times$ & \textbf{1.4$\times$} \\
    5-run total     & \$3.075  & \textbf{\$0.135} \\
    \bottomrule
  \end{tabular}
\end{table}

Table~\ref{tab:ch4_nd2_fix} shows the before/after comparison.  Patrol rate rises from
0\% to 100\%; quality rises from 0 to 4.0/5; cost per run falls from \$0.615 to
\$0.027 (96\% reduction).  The context growth drops from 34.2$\times$ to 1.4$\times$ because
M1 suppresses advisory re-injection and M2 replaces multi-thousand-token tool-result dumps
with compact structured state summaries ($\approx$240~tokens/turn).

\begin{figure}[ht]
  \centering
  \includegraphics[width=0.85\linewidth]{figures/fig4_5_nd2_fix_comparison.pdf}
  \caption{M1--M6 mitigation effect on GPT-4o-mini alert\_room (ND2).
    Left panel: patrol rate, quality, and relative cost (normalised to pre-fix baseline).
    Right panel: per-turn context size before and after mitigations.}
  \label{fig:ch4_nd2_fix}
\end{figure}

\subsection{ND2 Summary}
\label{sec:ch4_nd2_summary}

ND2 answers RQ2 (\textbf{yes}: repeated advisory injection produces LITM-class context
degradation) and provides the M1--M6 framework as a validated mitigation.  GPT-4o emerges
as the most balanced orchestrator for adversarial patrol scenarios.  GPT-4o-mini is fully
capable in benign conditions but requires loop mitigations for alert scenarios.  Gemini is
the cost-efficient tier for bulk patrol.  Claude shows an unexpected failure in clear-room
reasoning, suggesting its success in Chapter~\ref{chap:3rd chap} depends on the clarity of
the initial command rather than multi-turn reasoning endurance.

% ─────────────────────────────────────────────────────────────────────────────
\section{ND3: Human-in-the-Loop Integration}
\label{sec:ch4_nd3}
% ─────────────────────────────────────────────────────────────────────────────

\subsection{HITL Gate Design and Motivation}
\label{sec:ch4_nd3_design}

ND3 introduces the \texttt{approve\_callback} described in
Section~\ref{sec:ch4_proposed_hitl}.  The HITL gate tests whether agentic orchestrators
that successfully operated autonomously in ND2 remain coherent and cost-efficient when every
actuator command requires human sign-off.  The gate also tests whether the approval
mechanism provides sufficient loop-breaking force to replace the M1--M6 mitigations --- an
important question for practitioners who may prefer operator intervention over prompt
engineering.

ND3 uses the same mission design as ND2: 4 orchestrators $\times$ 2 scenarios $\times$
5 runs $= 40$ initial trials.  The ND3 mission prompt is \emph{identical} to ND2 (no
HITL-specific language added), testing whether the model adapts to the approval gate
through the denial feedback alone.

\subsection{Bug Discovery and Fixes}
\label{sec:ch4_nd3_bugs}

ND3 exposed two software bugs in the orchestrator loop, invalidating portions of the
initial results.  Both were fixed before re-running.

\paragraph{Bug A --- Gemini \texttt{finishReason="STOP"} break condition (Critical).}
The original code broke the agentic loop when \texttt{finish\_reason $\in$ \{STOP, MAX\_TOKENS,
SAFETY, \ldots\}} \emph{or} there were no function calls.  The Google Gemini API returns
\texttt{finishReason="STOP"} alongside function calls as its normal function-call
completion signal --- unlike OpenAI (\texttt{finish\_reason="tool\_calls"}) and Claude
(\texttt{stop\_reason="end\_turn"} only on text
turns)~\cite{geminiapi2025}.  The \texttt{or} conjunction caused the loop to break after a
single turn even when tool calls were present, yielding 0 drone commands, 0 approvals, and
0 cost across all 10 Gemini runs.

\textbf{Fix:} separate the conditions --- break only if there are no function calls OR
if \texttt{finish\_reason $\in$ \{MAX\_TOKENS, SAFETY, RECITATION, OTHER\}}; \texttt{STOP}
with function calls present is normal Gemini behaviour, continue.

\paragraph{Bug B --- Inner API error silently swallowed (Minor).}
An Azure 500 InternalServerError during GPT-4o-mini runs caused the inner retry loop to
\texttt{break} silently, leaving \texttt{error\_msg = ""} in the run record.  This made
API-outage runs appear as successful zero-data completions.

\textbf{Fix:} replace \texttt{break} with \texttt{raise} so exceptions propagate to the
run-level handler and are correctly recorded.

These bugs illustrate the importance of API-level testing when building multi-model
orchestration systems: provider conventions for finish reasons differ in subtle but
consequential ways.

\subsection{ND3B: Bug-Fixed Results}
\label{sec:ch4_nd3b}

After applying Bugs A and B fixes, three sub-experiments were run: ND3B (Gemini re-run,
10~runs), ND3B-mini (GPT-4o-mini re-run with Bug~B fix, 10~runs).  Claude and GPT-4o
results from the original ND3 were not affected by either bug and remain valid.

\begin{table}[ht]
  \caption{ND3 consolidated bug-fixed results (all orchestrators, both scenarios).
    GPT-4o alert\_room patrol 20\% reflects deliberate hazard-abort behaviour (correct
    safety doctrine).  GPT-4o-mini alert\_room 0\% is the LITM semantic loop;
    this is resolved in ND3C.  Gemini costs are 70$\times$ cheaper than Claude per run.}
  \label{tab:ch4_nd3b}
  \centering
  \small
  \begin{tabular}{llcccccc}
    \toprule
    Orchestrator & Scenario & Patrol\% & Quality/5 & Cost/Run & Turns & Approvals & Denials \\
    \midrule
    \multirow{2}{*}{Claude 3.5 Sonnet}
      & clear\_room & 60\%   & 3.60 & \$1.69 & 49.2 & 14.8 & 0 \\
      & alert\_room & 80\%   & 2.80 & \$2.13 & 39.4 & 14.4 & 0 \\
    \midrule
    \multirow{2}{*}{GPT-4o}
      & clear\_room & \textbf{100\%} & \textbf{4.00} & \$0.28 & 11.0 & 7.0 & 0 \\
      & alert\_room & 20\%$^\dagger$ & \textbf{4.20} & \$0.37 & 8.8  & 3.0 & 0 \\
    \midrule
    \multirow{2}{*}{Gemini 2.5 Flash (ND3B)}
      & clear\_room & 80\%   & 3.00 & \$0.023 & 38.8 & 10.6 & 0 \\
      & alert\_room & 60\%   & 4.00 & \$0.031 & 29.2 & 9.6  & 0 \\
    \midrule
    \multirow{2}{*}{GPT-4o-mini (ND3B)}
      & clear\_room & \textbf{100\%} & \textbf{4.00} & \$0.065 & 31.0 & 12.0 & 0 \\
      & alert\_room & 0\%$^\ddagger$ & 0.00 & \$1.70 & 50.0 & 13.0 & 10 \\
    \bottomrule
    \multicolumn{8}{l}{\footnotesize $\dagger$ Hazard-abort behaviour: model correctly
      prioritised emergency escalation over patrol completion.} \\
    \multicolumn{8}{l}{\footnotesize $\ddagger$ LITM semantic loop; resolved in ND3C
      (Table~\ref{tab:ch4_nd3c}).}
  \end{tabular}
\end{table}

\paragraph{Finding ND3-1 --- GPT-4o best quality, lowest cost among valid models.}
GPT-4o delivers quality~4.0 (clear\_room) and 4.2 (alert\_room) at \$0.28--\$0.37/run.
Clear\_room 100\% patrol is achieved in 11~turns on average --- the most efficient
full-patrol mission across all orchestrators in ND3.  The 20\% alert\_room patrol rate
reflects GPT-4o's hazard-abort priority: 4 of 5 runs write an emergency report and
request RTH immediately upon detecting the person, without completing all waypoints.  This
is correct real-world drone safety doctrine, supported by the SAFE-AI
framework~\cite{sanneman2022safeai}.

\paragraph{Finding ND3-2 --- Claude HITL-compliant but verbose and expensive.}
Claude received 100\% approval on all proposed commands (146/146, 0 denials) across both
scenarios, confirming it generates only safety-compliant tool calls.  However, at \$1.69
and \$2.13 per run respectively and 39--49~turns per run, Claude is 7$\times$ more expensive
than GPT-4o for comparable or lower quality.  The 60\% clear\_room patrol rate reflects two
runs where Claude saturated the 50-turn budget in a thorough sequential analysis without
writing the final report --- a mild LITM effect at high turn counts.

\paragraph{Finding ND3-3 --- Gemini 70$\times$ cheaper than Claude with near-comparable quality.}
At \$0.023--\$0.031/run and 70--80\% patrol rates, Gemini 2.5 Flash delivers
cost-per-mission economics suitable for high-frequency autonomous patrol at scale.
One run (clear\_room run~3) produced an anomaly: Gemini restarted the entire startup
sequence after observing an intermediate altitude reading ($0.781$~m $<$ 1.0~m target),
treating convergence-in-progress as a startup failure.  This probabilistic fixation on an
intermediate sensor value caused the run to complete only one waypoint.  The other 9 runs
operated normally.

Gemini also exhibited risk over-escalation: in clear\_room, 2/5~runs reported
\texttt{ALERT}/\texttt{EMERGENCY} for a room that MediaPipe scored as \texttt{CLEAR}.  In
alert\_room, 3/5~runs reported \texttt{EMERGENCY} for a \texttt{ALERT} ground truth.
Risk calibration via few-shot prompt examples is a viable mitigation for future work.

\begin{figure}[ht]
  \centering
  \includegraphics[width=0.85\linewidth]{figures/fig4_7_nd3b_results.pdf}
  \caption{ND3 consolidated results (bug-fixed).  Left: patrol rate by orchestrator and
    scenario.  Right: quality/5 vs.\ cost/run.  GPT-4o quality-cost frontier; Gemini
    cost-efficiency frontier.  GPT-4o-mini alert\_room excluded pending ND3C fix.}
  \label{fig:ch4_nd3b_results}
\end{figure}

\subsection{LITM Persists Under HITL: ND3B-mini Analysis}
\label{sec:ch4_nd3_mini}

A critical finding of ND3B-mini is that the HITL approval gate, while correctly denying
excessive hovering, is \emph{insufficient to break the semantic loop} on its own.

In GPT-4o-mini alert\_room (ND3B, no fixes), the approval gate denied \texttt{hover} on
every 4th+ call with the message ``\textit{excessive hovering --- mission requires forward
progress}''.  The model responded by substituting \texttt{navigate\_to\_waypoint} --- then
immediately proposing \texttt{hover} again.  The resulting deterministic
hover$\rightarrow$DENY$\rightarrow$navigate$\rightarrow$hover cycle ran until turn~50 on
all 5 runs, consuming \$1.70/run (\$8.50 for 5 runs) with 0 patrol completions.

Context grew 34.2$\times$ (12,871~tokens at turn~1 to 440,057~tokens at turn~50), matching
the standalone ND2 LITM trajectory exactly.  The HITL gate redirected the loop's immediate
action but provided no \emph{exit condition} the model could act on autonomously ---
answering RQ3 provisionally: operator denial alone is \emph{not} sufficient; semantic
loop mitigations are necessary even when a HITL gate is present.

This result extends the findings of~\citet{reflexion2023} on verbal reinforcement learning:
a correction signal (``stop hovering'') changes the immediate behaviour but not the
underlying policy unless the model has an internal rule mapping ``denied $n$ times
$\rightarrow$ change strategy''.  M4 and M6 provide exactly that rule.

\subsection{ND3C: M1--M6 + HITL --- Final Fix}
\label{sec:ch4_nd3c}

ND3C applies all six mitigations (Table~\ref{tab:ch4_nd2_mitigations}) to the GPT-4o-mini
alert\_room HITL scenario.  One ND3-specific adaptation: M4 tracks \texttt{hover\_count}
as the primary loop detector (rather than \texttt{analyze\_scene\_count} used in ND2),
and the nudge message explicitly references the operator denial pattern:
\textit{``You have been hovering repeatedly.  The operator has denied hover.  Write your
final report NOW, then call land().''}  M6's EXIT RULE likewise triggers on
\texttt{hover $\geq 3\times$} in addition to \texttt{analyze\_scene $\geq 3\times$}.

\begin{table}[ht]
  \caption{ND3C results: M1--M6 + HITL applied to GPT-4o-mini alert\_room (5 runs).
    All five runs complete the full patrol and write a report.  Zero hover denials ---
    the model self-terminates before the loop can form.  Quality 4.60/5 exceeds the ND2
    fix quality (4.0/5), showing that HITL approval acts as a secondary quality assurance
    layer.}
  \label{tab:ch4_nd3c}
  \centering
  \small
  \begin{tabular}{ccccccc}
    \toprule
    Run & Patrol & Quality/5 & Turns & Cost & Approvals & Denials \\
    \midrule
    1 & \checkmark & 3 & 8  & \$0.0148 & 3 & 0 \\
    2 & \checkmark & 5 & 9  & \$0.0165 & 2 & 0 \\
    3 & \checkmark & 5 & 8  & \$0.0148 & 3 & 0 \\
    4 & \checkmark & 5 & 9  & \$0.0165 & 3 & 0 \\
    5 & \checkmark & 5 & 9  & \$0.0165 & 3 & 0 \\
    \midrule
    \textbf{Mean} & \textbf{100\%} & \textbf{4.60} & \textbf{8.6} & \textbf{\$0.0158} & \textbf{2.8} & \textbf{0} \\
    \bottomrule
  \end{tabular}
\end{table}

\begin{table}[ht]
  \caption{ND3C before/after comparison: GPT-4o-mini alert\_room.  M1--M6 achieve a
    99\% cost reduction and restore 100\% patrol rate.  Context growth drops from 34.2$\times$
    to 1.45$\times$ --- an 8.6$\times$ reduction in per-run token volume.}
  \label{tab:ch4_nd3c_compare}
  \centering
  \small
  \begin{tabular}{lccc}
    \toprule
    Metric & ND3B (no fixes) & ND3C (M1--M6+HITL) & Change \\
    \midrule
    Patrol rate      & 0\%      & \textbf{100\%}   & $+100$~pp \\
    Quality/5        & 0.00     & \textbf{4.60}    & $+4.60$ \\
    Words written    & 0        & \textbf{128}     & $+128$ \\
    Turns/run        & 50~(max) & \textbf{8.6}     & $-83\%$ \\
    Cost/run         & \$1.70   & \textbf{\$0.016} & $-99\%$ \\
    Context growth   & 34.2$\times$ & \textbf{1.45$\times$} & $-96\%$ \\
    Hover denials/run & 10      & \textbf{0}       & $-100\%$ \\
    5-run total      & \$8.50   & \textbf{\$0.079} & $-99.1\%$ \\
    \bottomrule
  \end{tabular}
\end{table}

Tables~\ref{tab:ch4_nd3c} and~\ref{tab:ch4_nd3c_compare} present the ND3C results.
The M1--M6 package eliminates all hover denials: the model self-terminates after 3~hover
calls, writes the report, and lands in 8--9~turns.  The loop never forms --- the HITL gate
has nothing to reject.  This provides the final answer to RQ3: \textbf{yes}, M1--M6 resolve
semantic loop failures in HITL contexts, and the HITL gate combined with M1--M6 yields
higher quality (4.60/5) than M1--M6 alone in ND2 (4.0/5), because operator-approved
actions carry implicit validation that the model uses to anchor its report.

The three-stage progression (broken $\rightarrow$ loop $\rightarrow$ fixed) is illustrated
in Figure~\ref{fig:ch4_nd3c_prog} and mirrors the ND2 progression exactly.

\begin{figure}[ht]
  \centering
  \includegraphics[width=0.85\linewidth]{figures/fig4_8_nd3c_progression.pdf}
  \caption{Three-stage progression for GPT-4o-mini alert\_room: ND3 original (Azure outage,
    0 turns), ND3B (LITM loop, 50 turns, \$1.70), ND3C (M1--M6+HITL, 8.6 turns, \$0.016).
    Each stage required one targeted intervention: Bug~B fix, then M1--M6 mitigations.}
  \label{fig:ch4_nd3c_prog}
\end{figure}

\subsection{Emergency Monitor Behaviour Under HITL}
\label{sec:ch4_nd3_monitor}

The background emergency monitor fired reliably in all alert\_room runs, confirming
independent detection of the person hazard regardless of LLM loop state.  In ND3 the HITL
gate and monitor operate concurrently: a human denial of a drone command and an emergency
advisory can both be visible in the same turn's context.

The three valid orchestrators responded to this combined signal differently:

\begin{itemize}
  \item \textbf{GPT-4o}: immediately wrote an emergency report and requested RTH in 4/5
    alert\_room runs --- within 3--5 turns, before completing all waypoints.  Literature
    supports this as optimal when human safety is at risk~\cite{sanneman2022safeai}.

  \item \textbf{Gemini}: wrote an emergency report and landed in 3/5 runs; completed the
    patrol before reporting in the remaining 2/5.  Gemini 2.5 Flash exhibits
    over-escalation (ALERT $\rightarrow$ EMERGENCY) in risk labels, suggesting systematic
    conservatism in hazard assessment.

  \item \textbf{Claude}: continued the patrol incorporating the advisory into the report
    narrative rather than aborting --- the most task-focused but least safety-escalating
    response.
\end{itemize}

This variation reflects meaningful differences in how each model weights task completion
vs.\ safety escalation under combined HITL and emergency signals --- a finding with direct
implications for deployment configuration~\cite{kapoor2024aiagents}.

\subsection{HITL Operator Statistics}
\label{sec:ch4_nd3_hitl_stats}

\begin{table}[ht]
  \caption{ND3 HITL operator statistics (combined ND3 + ND3B valid runs).
    Three of four orchestrators operated 100\% within the auto-approver's safety bounds
    without a single denial.  Only GPT-4o-mini alert\_room triggered denials (hover loop).
    Mean operator response time (simulated) is 1.22--1.26~s per decision.}
  \label{tab:ch4_nd3_hitl}
  \centering
  \small
  \begin{tabular}{lcccc}
    \toprule
    Orchestrator & Approvals & Denials & Approval Rate & Mean Response (s) \\
    \midrule
    Claude 3.5 Sonnet & 146 & 0  & 100\% & 1.22 \\
    GPT-4o            & 50  & 0  & 100\% & 1.22 \\
    Gemini 2.5 Flash  & 101 & 0  & 100\% & 1.26 \\
    GPT-4o-mini (ND3B, alert\_room) & 65 & 50 & 57.7\% & 1.25 \\
    GPT-4o-mini (ND3C, M1--M6)     & 14 & 0  & \textbf{100\%} & 1.36 \\
    \bottomrule
  \end{tabular}
\end{table}

Table~\ref{tab:ch4_nd3_hitl} shows that three of four orchestrators never triggered a
single approval denial.  Claude, GPT-4o, and Gemini all generated tool calls that stayed
within safety bounds (\textpm 2.5~m patrol zone, 0.5--2.0~m altitude, $\leq$3~hover
calls, $\leq$2.0~m move step) without any explicit HITL-aware instruction in the prompt.
This suggests these models internalise conservative navigation strategies sufficient to
satisfy a realistic safety envelope.

The simulated mean operator response time of 1.22--1.36~s per decision adds approximately
12~s of latency over a 10-approval mission --- negligible on indoor surveillance timescales
(missions ran 24--200~s wall time).  For real deployments, UI-based approval (button tap)
would further reduce this latency~\cite{chen2018humanperformance}.

\subsection{ND3 Summary}
\label{sec:ch4_nd3_summary}

ND3 answers RQ3 (\textbf{yes, M1--M6 are necessary even with HITL; operator denial alone
cannot break a semantic loop}) and RQ4 (\textbf{GPT-4o best quality at reasonable cost;
Gemini best cost; Claude HITL-compliant but expensive}).  The HITL approval gate acts as an
additional quality assurance mechanism when combined with M1--M6, raising GPT-4o-mini
quality from 4.0 (ND2 fix) to 4.60 (ND3C).

% ─────────────────────────────────────────────────────────────────────────────
\section{Hardware Validation --- ND Series Live Drone Patrol (ND-HW)}
\label{sec:ch4_ndhw}
% ─────────────────────────────────────────────────────────────────────────────

\textit{This section is reserved for hardware validation results from a live drone patrol
experiment using the custom quadrotor platform described in Chapter~\ref{chap:3rd chap}
(Section~\ref{sec:ch3_b1hw}).  The experiment (\texttt{exp\_ND\_HW\_patrol.py}) will
replicate the ND3 clear\_room patrol scenario on the physical drone with the GPT-4o
orchestrator and the HITL gate operating via a mobile phone operator interface.
Results (patrol completion, HITL approval latency, risk label accuracy, cost) will be
reported here once data are available.}

\paragraph{Expected differences from simulation.}
Real-world ND-HW will introduce factors not present in simulation: (1) camera latency
($\approx$100~ms Raspberry Pi capture-to-USB pipeline) adds to \texttt{analyze\_scene}
response time; (2) MediaPipe false-positive rate on a real camera stream will differ from
the static test images used in ND1--ND3; (3) physical waypoint tracking errors
(expected: $\pm$5--10~cm based on B2 disturbance rejection results in
Chapter~\ref{chap:3rd chap}) will produce slightly different sensor telemetry fed back to
the LLM.

\paragraph{Success criteria.}
The live patrol will be considered successful if: patrol completion rate $\geq$80\%,
risk label accuracy $\geq$80\%, HITL approval latency $<$5~s per decision, and no
safety-bound violation occurs across a 5-run series.

% ─────────────────────────────────────────────────────────────────────────────
\section{Chapter Summary}
\label{sec:ch4_summary}
% ─────────────────────────────────────────────────────────────────────────────

This chapter presented the ND series, which closes the perception-action loop by combining
the Chapter~\ref{chap:2nd chap} image verbalization capability with the Chapter~\ref{chap:3rd chap}
tool-call execution framework into a unified agentic patrol system.

\paragraph{Key findings.}

\begin{enumerate}
  \item \textbf{Camera-as-tool is viable (ND1).}  All four LLM orchestrators call
    \texttt{analyze\_scene} autonomously at 100\% invocation rate.  GPT-4o-mini leads on
    quality (4.42/5); Gemini leads on cost (\$0.00043/run, 146$\times$ cheaper than Claude).

  \item \textbf{LITM is a real failure mode in agentic patrol (ND2).}  Repeated emergency
    monitor advisories cause GPT-4o-mini to enter a deterministic semantic loop in
    alert\_room scenarios, consuming 34.2$\times$ context growth and failing to write a
    single patrol report across 5 runs.  This confirms the U-shaped attention degradation
    of~\citet{liu2023lostinthemiddle} in an agentic drone control setting.

  \item \textbf{M1--M6 mitigations resolve LITM (ND2 fix + ND3C).}  The six-mitigation
    package transforms 0\% patrol / \$0.615 per run into 100\% patrol / \$0.027 per run in
    ND2, and 0\% / \$1.70 into 100\% / \$0.016 in the ND3C HITL context.  Context growth
    drops from 34.2$\times$ to 1.45$\times$.

  \item \textbf{HITL alone cannot break semantic loops (ND3B-mini).}  The approval gate
    correctly identifies excessive hovering and denies it, but denial redirects rather than
    terminates the loop.  M4+M6 are necessary to give the model an internal exit condition.

  \item \textbf{HITL + M1--M6 achieves highest quality (ND3C).}  GPT-4o-mini with all
    fixes and the HITL gate achieves 4.60/5 quality --- the highest across all ND series
    experiments.  Operator-approved actions provide implicit validation that the model
    leverages in report writing.

  \item \textbf{Two API bugs discovered (ND3).}  The Gemini \texttt{finishReason="STOP"}
    break condition and the OpenAI/Gemini silent error swallow were found and fixed,
    restoring Gemini from 1-turn (0 commands) to full multi-waypoint missions.
\end{enumerate}

\paragraph{Production configuration.}

\begin{table}[ht]
  \caption{ND series recommended production configuration based on cost-quality tradeoffs.
    Three deployment profiles: quality-maximising (GPT-4o), cost-efficient (Gemini),
    and budget HITL (GPT-4o-mini + M1--M6).}
  \label{tab:ch4_production}
  \centering
  \small
  \renewcommand{\arraystretch}{1.2}
  \begin{tabular}{p{2.8cm}p{2.5cm}p{1.6cm}p{1.6cm}p{4.5cm}}
    \toprule
    Profile & Orchestrator & Quality/5 & Cost/Run & Notes \\
    \midrule
    Quality-maximising & GPT-4o & 4.10 & \$0.32 &
      Best report quality; correct hazard-abort safety doctrine \\
    Cost-efficient & Gemini 2.5 Flash & 3.50 & \$0.027 &
      70$\times$ cheaper than Claude; suitable for bulk patrol;
      finishReason bug must be patched \\
    Budget HITL & GPT-4o-mini + M1--M6 & 4.60 & \$0.016 &
      Highest quality in HITL setting; lowest cost; M1--M6 required for
      alert scenarios \\
    \midrule
    Not recommended & Claude 3.5 Sonnet & 3.20 & \$1.91 &
      HITL-compliant but 7$\times$ more expensive than GPT-4o for lower quality \\
    \bottomrule
  \end{tabular}
\end{table}

\begin{figure}[ht]
  \centering
  \includegraphics[width=0.80\linewidth]{figures/fig4_9_production_tradeoff.pdf}
  \caption{ND series production tradeoff: quality/5 vs.\ cost/run.  The quality-cost
    Pareto frontier runs from Gemini 2.5 Flash (bottom-right, cheapest) through
    GPT-4o-mini+M1--M6 (high quality, low cost) to GPT-4o (highest quality, moderate cost).
    Claude lies off the frontier (high cost, lower quality).}
  \label{fig:ch4_production_tradeoff}
\end{figure}

\paragraph{Connection to Chapter~\ref{chap:5th chap}.}
The ND series establishes that the three-tier architecture (PID $\rightarrow$ vision $\rightarrow$
LLM orchestrator) operates correctly in simulation across a range of orchestrators, failure
modes, and HITL configurations.  The next chapter synthesises the findings across
Chapters~\ref{chap:2nd chap}--\ref{chap:4th chap}, discusses limitations (simulation
fidelity, single-indoor-scenario scope, absence of real-time closed-loop telemetry during
LLM reasoning), and proposes future work directions including adaptive HITL threshold
calibration, multi-room patrol graph extension, and on-device LLM deployment for
latency-critical applications.
